\section{Problem formulation}
\label{sec:problem}

\subsection{Trajectory state and belief}

A scientific-agent trajectory at step $t$ produces a structured event
history $\mathcal{H}_t = \{e_0, \dots, e_t\}$ from a typed event schema
that distinguishes proposals, action requests, action results, state
snapshots, state differences, evidence, verifier outcomes, and fault
injection records~\citep{reliability_benchmark}. The frozen R1 benchmark supplies a
calibrated belief state:

\begin{align}
    b_t &= (q_t, \mathbf{p}_t, f_t, r_t, d_t) \label{eq:belief}
\end{align}

where $q_t = P(\text{terminal scientific failure} \mid \mathcal{H}_t)$ is
the estimated failure hazard, $\mathbf{p}_t$ is a distribution over causal
fault families, $f_t \in [0,1]$ is evidence freshness, $r_t \in [0,1]$ is
estimated reversibility (probability that a recovery action can restore a
valid state), and $d_t \in [0,1]$ is estimated accumulated damage. The
belief is fixed at the version supplied by the R1 artifact; R2 does not
retrain the belief estimator.

\subsection{Action space and cost model}

At each eligible decision point, the controller selects from:

\begin{align}
    a_t \in \{\, &\texttt{continue},\,
                 \texttt{verify}(c),\,\nonumber\\
                 &\texttt{recover}(\texttt{retry} \mid \texttt{repair} \mid
                                   \texttt{rollback} \mid \texttt{re-plan}),\,\nonumber\\
                 &\texttt{escalate},\,
                 \texttt{abstain},\,
                 \texttt{terminate} \,\} \label{eq:actions}
\end{align}

Each action carries an explicit cost vector $\mathbf{c}(a_t) =
(c_v, c_{\text{compute}}, c_{\text{human}}, c_{\text{damage}})$, where
$c_v(a_t)$ is the verification or recovery execution cost, $c_{\text{compute}}$
is tool-call and compute cost, $c_{\text{human}}$ is human escalation cost,
and $c_{\text{damage}}$ is the estimated risk of irreversible damage.
The cost vector is declared for each action before evaluation; sensitivity
to cost-model choices is tested as a robustness analysis.

A verification action $\texttt{verify}(c)$ selects a check type $c$ and a
target component. The expected benefit of verification is the reduction in
uncertainty about failure status or cause, weighted by the cost of acting
on an incorrect belief. A recovery action specifies the recovery subclass
and includes declared preconditions (reversibility within budget, known
rollback boundary) and postconditions (fault removed, downstream state
consistent, no new validity violations).

\subsection{Budget-constrained utility}

The controller optimizes expected utility subject to budget and safety
constraints. Let $\mathcal{T}$ be the set of trajectories. The objective is:

\begin{align}
    \max_{\pi} \; \mathbb{E}_{\tau \sim \pi}\Big[&
        U(\tau_{\text{outcome}})
        - \lambda_v \cdot C_v(\tau)
        - \lambda_{\text{compute}} \cdot C_{\text{compute}}(\tau) \nonumber\\
        &- \lambda_{\text{human}} \cdot C_{\text{human}}(\tau)
        - \lambda_{\text{damage}} \cdot D(\tau) \Big] \\
    \text{s.t.} \quad & \mathbb{E}_{\tau_{\text{clean}} \sim \pi}[U_{\text{clean}}]
        \geq U_{\text{base}} - \delta_{\text{regression}} \label{eq:regression}\\
    & \text{selective risk} \leq R_{\text{budget}} \nonumber
\end{align}

where $U(\tau_{\text{outcome}})$ is the adjudicated scientific utility,
$C_v$, $C_{\text{compute}}$, $C_{\text{human}}$ are cumulative costs,
$D(\tau)$ is irreversible damage, and $\lambda$ coefficients trade off
costs against utility. The clean-task regression constraint
(Eq.~\ref{eq:regression}) requires that the policy does not degrade
performance on tasks that would have succeeded without intervention beyond
a predeclared tolerance $\delta_{\text{regression}} = 0.02$ in the primary
analysis.

\subsection{Causal interpretation of randomized controller effects}

The randomized online study (R2-C1) assigns policies within blocks defined
by task instance, model configuration, and sealed fault episode. Under
block randomization with the standard stable-unit-treatment-value
assumption (each trajectory's outcome depends only on its assigned policy,
not on policies assigned to other trajectories in the same block), the
difference in independently verified scientific success between policies
has a causal interpretation as the average treatment effect of the policy
within the sampled population of tasks, faults, and models.

\subsection{Assumptions}

We state the following assumptions explicitly. (1)~The R1 belief $b_t$ is a
sufficient statistic for the information in $\mathcal{H}_t$ relevant to
failure prediction; the controller does not have access to hidden fault
truth. (2)~Adjudication is independent of the controller and the training
reward channel. (3)~Training and evaluation verifiers are separated where
feasible; when they must share an implementation, this is declared in the
reward-channel audit. (4)~No task, fault implementation, near-duplicate
trace, or adjudicator-derived label crosses the training/evaluation
boundary. (5)~Unsupported counterfactual actions (those not observed under
any policy with adequate propensity support) are treated as explicitly
unknown rather than imputed.
